\chapter*{\IfLanguageName{dutch}{Woord vooraf}{Preface}}%
\label{ch:voorwoord}

Deze bachelorproef vormt het sluitstuk van mijn opleiding Toegepaste Informatica, afstudeerrichting AI Data Engineering. Met dit onderzoek naar datareductie binnen AI-visionmodellen heb ik bewust gekozen voor een onderwerp dat technische diepgang combineert met maatschappelijke relevantie. De voortdurende schaalvergroting binnen deep learning – gekenmerkt door steeds grotere datasets en toenemende rekenkracht – roept fundamentele vragen op over efficiëntie, duurzaamheid en verantwoord gebruik van middelen. Vanuit die interesse is dit onderzoek gegroeid.

Tijdens het uitwerken van deze bachelorproef heb ik mijn kennis rond machine learning en computer vision verder kunnen verdiepen, maar ook geleerd om wetenschappelijke literatuur kritisch te analyseren en te vertalen naar een concrete experimentele methodologie. Het opzetten van reproduceerbare experimenten, het objectief evalueren van resultaten en het afwegen van performantie tegenover energieverbruik en trainingstijd vormden een bijzonder leerrijk traject.

Mijn oprechte dank gaat uit naar mijn promotor, Johan De Corte, voor zijn academische begeleiding, kritische reflecties en waardevolle feedback gedurende het volledige proces. Zijn inzichten hebben mij geholpen om mijn onderzoeksvraag scherp te stellen en mijn werk inhoudelijk te versterken.

Daarnaast wens ik mijn co-promotor, Arshia Shariatmadar, te bedanken voor zijn technische expertise en praktijkgerichte ondersteuning binnen het domein van artificiële intelligentie. Zijn professionele inzichten hebben bijgedragen aan de realiteitszin en toepasbaarheid van dit onderzoek.

Tot slot wil ik mijn familie en naasten bedanken voor hun voortdurende steun en vertrouwen tijdens mijn opleiding en het schrijven van deze bachelorproef. Hun aanmoediging heeft mij geholpen dit traject met doorzettingsvermogen af te ronden.